\chapter{場所と道案内 — Lugares y Direcciones}
\unidadheader{04}{場所と道案内}{Lugares y Direcciones}{Preguntar y dar indicaciones}

\begin{tcolorbox}[vocabbox, title={📌 Vocabulario Nuevo}]
\begin{tabularx}{\linewidth}{llX}
\toprule
\textbf{Japonés} & \textbf{Español} & \textbf{Tipo} \\
\midrule
\vocabitem{駅}{えき}{estación (tren/metro)}{sust.}
\vocabitem{学校}{がっこう}{escuela}{sust.}
\vocabitem{病院}{びょういん}{hospital}{sust.}
\vocabitem{銀行}{ぎんこう}{banco}{sust.}
\vocabitem{図書館}{としょかん}{biblioteca}{sust.}
\vocabitem{右}{みぎ}{derecha}{sust.}
\vocabitem{左}{ひだり}{izquierda}{sust.}
\vocabitem{前}{まえ}{frente / adelante}{sust.}
\vocabitem{後ろ}{うしろ}{atrás}{sust.}
\vocabitem{近く}{ちかく}{cerca (de)}{sust./adv.}
\bottomrule
\end{tabularx}
\end{tcolorbox}

\subsection*{🗣️ Frases Clave}

\frase{〜はどこですか？}{¿Dónde está ~?}
\frase{〜の\ruby{近}{ちか}くにあります。}{Está cerca de ~.}
\frase{まっすぐ\ruby{行}{い}ってください。}{Vaya recto, por favor.}
\frase{\ruby{右}{みぎ}に\ruby{曲}{ま}がってください。}{Doble a la derecha, por favor.}
\frase{\ruby{二}{に}つ\ruby{目}{め}の\ruby{角}{かど}です。}{Es la segunda esquina.}

\subsection*{⚙️ Gramática}

\patron{Hay / Existe (cosas y lugares)}
  {〜が あります　／　〜がいます}
  {あります = cosas inanimadas \quad | \quad います = seres animados
  \ej{\ruby{駅}{えき}の\ruby{近}{ちか}くにコンビニがあります。}
    {Hay un convenience store cerca de la estación.}
  \ej{そこに\ruby{先生}{せんせい}がいます。}{Hay un maestro ahí.}}

\patron{Partículas de lugar}
  {〜に (existencia)　／　〜で (acción)　／　〜へ・に (dirección)}
  {\ej{\ruby{図書館}{としょかん}に\ruby{本}{ほん}があります。}
    {Hay libros en la biblioteca. (existencia)}
  \ej{\ruby{図書館}{としょかん}で\ruby{勉強}{べんきょう}します。}
    {Estudio en la biblioteca. (acción)}
  \ej{\ruby{学校}{がっこう}へ\ruby{行}{い}きます。}
    {Voy a la escuela. (dirección)}}

\begin{tcolorbox}[ejerciciobox, title={📝 Ejercicios Rápidos}]
\begin{enumerate}
  \item ¿に o で? Completa: コンビニ\_\_\_ジュースを\ruby{買}{か}います。
  \item Pregunta dónde está el hospital en japonés.
  \item Traduce: \ruby{駅}{えき}の\ruby{右}{みぎ}に\ruby{銀行}{ぎんこう}があります。
\end{enumerate}
\vspace{0.4em}
\textbf{Respuestas:}
1) で\quad
2) \ruby{病院}{びょういん}はどこですか？\quad
3) Hay un banco a la derecha de la estación.
\end{tcolorbox}

\begin{tcolorbox}[kanjibox, title={🀄 Kanjis de la Unidad}]
\begin{tabularx}{\linewidth}{cllXX}
\toprule
\textbf{Kanji} & \textbf{On} & \textbf{Kun} & \textbf{Significado} & \textbf{Vocab} \\
\midrule
\kanjirow{駅}{えき}{—}{estación}{\ruby{駅}{えき} / \ruby{終点}{しゅうてん}}
\kanjirow{右}{ゆう}{みぎ}{derecha}{\ruby{右}{みぎ} / \ruby{右手}{みぎて}}
\kanjirow{左}{さ}{ひだり}{izquierda}{\ruby{左}{ひだり} / \ruby{左手}{ひだりて}}
\kanjirow{前}{ぜん}{まえ}{frente / antes}{\ruby{前}{まえ} / \ruby{前後}{ぜんご}}
\kanjirow{近}{きん}{ちか}{cerca}{\ruby{近く}{ちかく} / \ruby{近所}{きんじょ}}
\bottomrule
\end{tabularx}
\end{tcolorbox}
